\documentclass[11pt]{article}
\usepackage[a4paper,margin=29mm]{geometry}
\usepackage[T1]{fontenc}
\usepackage{lmodern,amsmath,amssymb,amsthm,mathtools,microtype}
\usepackage[colorlinks=true,linkcolor=blue,citecolor=blue,urlcolor=blue]{hyperref}
\newtheorem{theorem}{Theorem}[section]
\newtheorem{lemma}[theorem]{Lemma}
\newtheorem{remark}[theorem]{Remark}
\newcommand{\R}{\mathbb R}
\newcommand{\eps}{\varepsilon}
\newcommand{\dd}{\,\mathrm d}
\newcommand{\Om}{\Omega}
\newcommand{\De}{D_\eps}
\newcommand{\qe}{q_\eps}
\newcommand{\norm}[1]{\left\lVert#1\right\rVert}
\newcommand{\av}[1]{\overline{#1}}
\DeclareMathOperator{\supp}{supp}
\title{The nonlocal-to-local limit for superlinear SKT cross-diffusion}
\author{Kaiwen Jin\\\small\href{mailto:yui.ui.siki@gmail.com}{yui.ui.siki@gmail.com}}
\date{October 15, 2026}
\begin{document}
\maketitle
\begin{abstract}
We prove convergence from regional nonlocal diffusion to the local
Shigesada--Kawasaki--Teramoto system with superlinear transition rates.
For any number of species and any fixed exponent greater than one, the given
nonnegative strong nonlocal solutions admit a subsequence converging strongly
in space-time on every finite time interval, under the symmetric coefficient
and strict diagonal-dominance assumptions of the source problem.
The limit satisfies the local equation with its original initial data and
homogeneous Neumann pressure flux. Both nonlinear pressures and reactions
converge strongly in $L^1$.
A power entropy gives compactness of the densities. A supplementary
regularized logarithmic entropy controls weighted square-root differences,
which yield an integrable pressure gradient and permit first-order testing
up to the boundary.
\end{abstract}

\section{Introduction and main result}\label{sec:intro}

Cross-diffusion describes populations whose dispersal depends on the densities
of other species. In a regional nonlocal model, individuals jump between points
of a bounded habitat, and the jump intensity depends on the population at the
departure and arrival points. When the jump length tends to zero, one expects
a local diffusion equation with no flux through the boundary. Establishing this
limit requires estimates uniform in the jump length; the strong bounds for an
individual nonlocal equation do not suffice.

For two species with linear transition rates, Galiano and Velasco
\cite[Theorem 1]{GV} proved such a localization result. Hirvonen, J\"ungel and
Pollino \cite[Theorems 1 and 2]{HJP} proved global existence and uniqueness of
strong solutions, together with entropy estimates, for a multispecies model
with superlinear transition rates. The
localization question for those rates is stated in \cite[Section 1.1]{HJP}
and in AIM problem 469 \cite{AIM}. Existence theorems for local nonlinear
cross-diffusion systems, such as \cite[Theorem 2]{CDJ}, do not identify limits
of the prescribed regional jump solutions. Moreover, density compactness alone
does not identify a pressure gradient or the first-order weak boundary
condition; the distinction from second-order testing is relevant here
\cite[Remark 16]{CDJ}.

We prove the localization result with all dimensions, species and exponents
in the source problem. The coefficient, kernel and initial-data hypotheses
below are exactly those of that problem. In addition to strong density
convergence, we obtain an integrable weak gradient of the actual nonlinear
pressure. This gives the homogeneous Neumann condition with unrestricted
boundary values of the test function.

\subsection{The equations and assumptions}

Let $d\geq1$, let $\Om\subset\R^d$ be a nonempty smooth bounded connected
domain, let $n\geq2$, and fix $s>1$. For $u=(u_1,\ldots,u_n)\geq0$, set
\begin{equation}\label{eq:PR}
 P_i(u)=u_i\left(a_{i0}+\sum_{j=1}^n a_{ij}u_j^s\right),
 \qquad R_i(u)=u_i\left(b_{i0}-\sum_{j=1}^n b_{ij}u_j\right).
\end{equation}
Assume $a_{i0}>0$, $b_{i0},b_{ij}\geq0$, and
$a_{ij}=a_{ji}\geq0$, with
\begin{equation}\label{eq:dominance}
 a_{ii}>C(s)\sum_{j\ne i}a_{ij},\qquad
 C(s)=\begin{cases}
 \displaystyle\frac{s^2}{2(s+1)(s-1)},&1<s<2,\\[4pt]
 \displaystyle\frac{4(s-1)^2+1}{4(s+1)(s-1)},&s\geq2.
 \end{cases}
\end{equation}
Let $J\in C_c^\infty(\R^d)$ be nonnegative, radial, and positive near zero,
with normalization
\[
 \frac12\int_{\R^d}J(z)z_1^2\dd z=1.
\]
Thus, by radiality,
\begin{equation}\label{eq:moment}
 \int_{\R^d}J(z)z_kz_l\dd z=2\delta_{kl},\qquad
 M_2:=\int_{\R^d}J(z)|z|^2\dd z=2d.
\end{equation}
Write $J_\eps(h)=\eps^{-d}J(h/\eps)$ and fix
$0\leq u_i^0\in L^\infty(\Om)\cap BV(\Om)$.
We consider the given global nonnegative strong solutions of
\begin{equation}\label{eq:nonlocal}
 \partial_tu_i^\eps(t,x)=\eps^{-2}\int_\Om J_\eps(x-y)
 \bigl[P_i(u^\eps(t,y))-P_i(u^\eps(t,x))\bigr]\dd y
 +R_i(u^\eps(t,x)),\quad u_i^\eps(0)=u_i^0.
\end{equation}
The strong sense is that, on each finite interval, the solution belongs to
$W^{1,\infty}(0,T;L^\infty(\Om))$ and
$C([0,T];L^\infty(\Om)\cap BV(\Om))$, and satisfies
\eqref{eq:nonlocal} almost everywhere. This is the strong-solution setting
of \cite[Theorem 1]{HJP}. No bound in these spaces uniform in $\eps$ is assumed.

\begin{theorem}\label{thm:main}
Under these assumptions, every sequence $\eps\to0$ has one subsequence and a
global nonnegative limit $u$ such that, for every $T<\infty$, writing
$Q_T=(0,T)\times\Om$ and
\[
 r=\frac{4s}{3s+1}>1,\qquad r'=\frac{4s}{s-1},
\]
we have
\begin{gather}
 u_i^\eps\longrightarrow u_i\quad\hbox{strongly in }L^p(Q_T)
       \quad(1\leq p<2s),\label{eq:conv}\\
 P_i(u^\eps)\longrightarrow P_i(u),\quad
 R_i(u^\eps)\longrightarrow R_i(u)\quad\hbox{strongly in }L^1(Q_T),
 \label{eq:productconv}\\
 \begin{aligned}
 u_i&\in L^\infty(0,T;L^s(\Om))\cap L^{2s}(Q_T)
                \cap C_w([0,T];L^s(\Om)),\\
 u_i^s,\sqrt{u_i}&\in L^2(0,T;H^1(\Om)),\qquad
 P_i(u)\in L^r(0,T;W^{1,r}(\Om)),\\
 R_i(u)&\in L^s(Q_T),\qquad
 \partial_tu_i\in L^r(0,T;(W^{1,r'}(\Om))').
 \end{aligned}\label{eq:regularity}
\end{gather}
Here $C_w$ denotes continuity for the weak $L^s$ topology. The limit has
$u_i(0)=u_i^0$ weakly in $L^s(\Om)$ and, for every
$\phi\in C^1([0,T]\times\overline\Om)$ with $\phi(T)=0$,
\begin{equation}\label{eq:weak}
 -\int_{Q_T}u_i\partial_t\phi-\int_\Om u_i^0\phi(0)
 +\int_{Q_T}\nabla P_i(u)\cdot\nabla\phi
 =\int_{Q_T}R_i(u)\phi.
\end{equation}
No restriction is imposed on the boundary values of $\phi$.
\end{theorem}

Equation \eqref{eq:weak} is the first-order weak formulation of
$\partial_tu_i=\Delta P_i(u)+R_i(u)$ with homogeneous Neumann pressure flux.
It does not assert a classical normal trace. The subsequence is chosen once
for all finite horizons. We do not prove uniqueness of the local weak solution,
a convergence rate, convergence of the whole family, or the endpoint $s=1$.

The proof has two complementary entropy estimates. The power entropy controls
differences of $u_i^s$ and yields higher integrability. To estimate differences
of the products $u_i u_j^s$ in the pressure, we also need information near
vacuum. A regularized logarithmic entropy supplies the corresponding weighted
square-root differences. Their combination gives a pressure difference
quotient in $L^r$, with $r>1$. A nonlinear compactness argument then gives
strong convergence, and the pressure quotient identifies the weak gradient
and allows testing through the boundary. Section~\ref{sec:entropy} establishes
the estimates; Section~\ref{sec:limit} proves compactness and the limit passage.

\section{Uniform entropy estimates}\label{sec:entropy}

All constants below may depend on the fixed domain, coefficients, $n,s,J,u^0$
and $T$, but never on $\eps$ or on a removed regularization parameter.
For a spatial function $w$, define the regional quadratic form
\begin{equation}\label{eq:D}
 \De(w)=\iint_{\Om^2}\eps^{-d-2}J((x-y)/\eps)
                    |w(y)-w(x)|^2\dd x\dd y.
\end{equation}
Its form domain is $L^2(\Om)$, and we set $\De(w)=+\infty$ outside that
space. Every fixed-$\eps$ solution slice used below is bounded, so its powers
are in this form domain. No uniform $L^2$ bound is implicit in the definition.

\subsection{Power entropy and spatial integrability}

Set $S_i=\sum_{j\ne i}a_{ij}$ and $k_i=(s+1)a_{ii}-(s-1)S_i$.
Since $(s+1)C(s)>s-1$, condition \eqref{eq:dominance} implies $k_i>0$.
For $1<s<2$ this follows from $s^2>2(s-1)^2$; for $s\geq2$ the difference
between $C(s)$ and $(s-1)/(s+1)$ is $1/[4(s+1)(s-1)]$.

\begin{lemma}\label{lem:power}
For nonnegative vectors $X,Y$, one has
\begin{align}
 \sum_i(X_i^{s-1}-Y_i^{s-1})(P_i(X)-P_i(Y))
 &\geq\frac{4(s-1)}{s^2}\sum_i a_{i0}|X_i^{s/2}-Y_i^{s/2}|^2\notag\\
 &\quad+\frac{s-1}{s^2}\sum_i k_i|X_i^s-Y_i^s|^2.\label{eq:coercivity}
\end{align}
Consequently the solutions satisfy
\begin{equation}\label{eq:powerbound}
 \sup_{0\leq t\leq T}\sum_i\norm{u_i^\eps(t)}_s^s
       +\sum_i\int_0^T\De((u_i^\eps)^s)\dd t\leq C_T.
\end{equation}
\end{lemma}
\begin{proof}
First take positive vectors. In the variables $v_i=u_i^{s-1}$, write
$F(v)=P(u)$. Direct differentiation gives
\[
 \partial_iF_i=\frac{a_{i0}u_i^{2-s}+(s+1)a_{ii}u_i^2
             +\sum_{j\ne i}a_{ij}u_i^{2-s}u_j^s}{s-1},\qquad
 \partial_jF_i=\frac{s}{s-1}a_{ij}u_iu_j\quad(i\ne j).
\]
This matrix is symmetric. With $z_i=u_i\xi_i$, its off-diagonal pair is
bounded using
\[
 a_{ij}\bigl[(u_j/u_i)^s z_i^2+(u_i/u_j)^s z_j^2+2sz_iz_j\bigr]
       \geq-(s-1)a_{ij}(z_i^2+z_j^2).
\]
Indeed, the first two terms are at least $2|z_iz_j|$, and the only adverse
case is $z_iz_j<0$. Hence
\begin{equation}\label{eq:Jacobian}
 \xi^TDF(v)\xi\geq\frac1{s-1}
          \sum_i(a_{i0}u_i^{2-s}+k_i u_i^2)\xi_i^2.
\end{equation}
Integrate along the straight segment from $Y^{s-1}$ to $X^{s-1}$, and put
$U_i(\theta)=((1-\theta)Y_i^{s-1}+\theta X_i^{s-1})^{1/(s-1)}$.
Differentiating $U_i^s$ and $U_i^{s/2}$ and applying Cauchy--Schwarz gives
\[
 (X_i^{s-1}-Y_i^{s-1})^2\int_0^1U_i^2\dd\theta
        \geq\frac{(s-1)^2}{s^2}|X_i^s-Y_i^s|^2,
\]
\[
 (X_i^{s-1}-Y_i^{s-1})^2\int_0^1U_i^{2-s}\dd\theta
        \geq\frac{4(s-1)^2}{s^2}|X_i^{s/2}-Y_i^{s/2}|^2.
\]
Together with \eqref{eq:Jacobian} these prove \eqref{eq:coercivity}.
Continuity extends that inequality to nonnegative endpoints. Thus no
fractional power is differentiated at vacuum.

Use $h_s(t)=(t^s-st+s-1)/(s-1)\geq0$ as entropy, and let
$E(t)=\sum_i\int_\Om h_s(u_i^\eps(t))$. Symmetry of the kernel gives
\begin{equation}\label{eq:entropy}
 E'(t)+\frac2s\sum_i a_{i0}\De((u_i^\eps)^{s/2})
      +\frac1{2s}\sum_i k_i\De((u_i^\eps)^s)
       \leq C(1+E(t)).
\end{equation}
For the reaction bound, if $u_i\geq1$, the competition term multiplied by
$h_s'(u_i)$ is nonpositive, and the growth term is bounded by $C(1+u_i^s)$.
If $u_i<1$, the positive competition contribution is bounded by
$C\sum_j u_j\leq C(1+\sum_j u_j^s)$. Finally
$u_i^s\leq C(1+h_s(u_i))$.

The chain rule is legitimate for every fixed $\eps$: the solution is a
bounded Lipschitz time curve into $L^\infty$, and the entropy is $C^1$ on
its nonnegative range. At zero use the $C^1$ extension
$(|t|^s-st+s-1)/(s-1)$. The derivative of $U\mapsto\int h_s(U)$ is
$V\mapsto\int h_s'(U)V$; the remainder is bounded by
$|\Om|\norm{V}_\infty$ times the modulus of continuity of $h_s'$ on the
fixed bounded range. Thus the usual absolutely continuous scalar chain rule
applies. All fixed-$\eps$ pair integrals are finite before symmetrization.
Multiplication of the differential inequality by $e^{-Ct}$ bounds $E$;
integration of \eqref{eq:entropy} then proves \eqref{eq:powerbound}.
\end{proof}

We use the following precise spatial consequence of Ponce's results. It
restores the mean before asserting a uniform norm bound.

\begin{lemma}\label{lem:Ponce}
For sufficiently small $\eps$ and every $w\in L^2(\Om)$,
\begin{equation}\label{eq:Poincare}
 \norm{w-\av w}_2^2\leq C\De(w),\qquad
 \av w=|\Om|^{-1}\int_\Om w.
\end{equation}
If $\eps_k\to0$ and both $\norm{w_k}_2$ and $D_{\eps_k}(w_k)$ are bounded,
then a subsequence converges strongly in $L^2(\Om)$.
\end{lemma}
\begin{proof}
Recall the second moment $M_2=\int J(z)|z|^2\dd z=2d$ and set
$\rho_\eps(h)=|h|^2J_\eps(h)/(M_2\eps^2)$.
These are nonnegative radial probability kernels concentrating at zero, and
\[
 \iint_{\Om^2}\frac{|w(y)-w(x)|^2}{|x-y|^2}\rho_\eps(x-y)\dd x\dd y
                 =\frac1{M_2}\De(w).
\]
For $d\geq2$, \eqref{eq:Poincare} and the compactness assertion are
\cite[Theorems 1.1--1.2]{Ponce}, at $p=2$.
Their hypotheses are a bounded Lipschitz domain, radial probability kernels
concentrating at zero, and the displayed bounded norm and energy.
The eventual Poincar\'e constant can be chosen uniformly for the present
family: otherwise choose a violating sequence $\eps_k\to0$ and apply
Theorem 1.1 to that very sequence with its fixed constant.

For $d=1$, connectedness makes $\Om$ an interval. Choose $R,j_0>0$ with
$J\geq j_0$ on $[-R,R]$. If $R\eps/2\leq |h|\leq R\eps$ and
$1/2\leq\theta\leq1$, then
\[
 \rho_\eps(\theta h)\geq\frac{j_0R^2}{16M_2\eps}.
\]
The annulus has length $R\eps$, so
$\int\inf_{1/2\leq\theta\leq1}\rho_\eps(\theta h)\dd h
\geq j_0R^3/(16M_2)>0$.
This verifies condition (8) of \cite[Theorem 1.3]{Ponce}; the kernels here
are smooth, so its Lebesgue-point representative causes no ambiguity.
That theorem extends both assertions to $(0,1)$. For
$\Om=(a,a+L)$, use $W(x)=w(a+Lx)$ and
$\widetilde\rho(h)=L\rho_\eps(Lh)$. The original weighted energy is
$L^{-1}$ times the energy on $(0,1)$, and the original squared oscillation
is $L$ times that on $(0,1)$. Only the fixed factor $L^2$ changes the
Poincar\'e constant. Compactness transfers by the same rescaling.
\end{proof}

Apply \eqref{eq:Poincare} to $w=(u_i^\eps)^s$. Its mean is bounded by
\eqref{eq:powerbound}. Thus, for all sufficiently small $\eps$,
\begin{equation}\label{eq:L2s}
 \sum_i\norm{u_i^\eps}_{L^{2s}(Q_T)}^{2s}\leq C_T.
\end{equation}
In particular $P_i(u^\eps)$ is bounded in $L^{2s/(s+1)}(Q_T)$ and
$R_i(u^\eps)$ in $L^s(Q_T)$. Both exponents exceed one.

\subsection{A logarithmic estimate at vacuum}

To control the pressure gradient, the power estimate must be supplemented by
differences of $\sqrt{u_i^\eps}$ with weights from the other species.
For time-dependent $w$, define
\[
 \qe w(t,x,z)=\mathbf1_\Om(x+\eps z)
              \frac{w(t,x+\eps z)-w(t,x)}{\eps},\qquad
 \dd\mu=\dd t\dd x\,J(z)\dd z.
\]
The quotient is zero when the translated point leaves $\Om$; it uses no
outside value of $w$. The measure $\mu$ on
$(0,T)\times\Om\times\R^d$ is finite and independent of $\eps$, and
$\norm{\qe w}_{L^2(\mu)}^2=\int_0^T\De(w(t))\dd t$.

\begin{lemma}\label{lem:log}
The solutions satisfy
\begin{align}
 \sum_i\int_0^T\De(\sqrt{u_i^\eps})\dd t
 +\sum_{i,j}a_{ij}\int
 &\min\{(u_j^\eps(t,x))^s,(u_j^\eps(t,x+\eps z))^s\}\notag\\[-2pt]
 &\hspace{8mm}\cdot|\qe\sqrt{u_i^\eps}|^2\dd\mu\leq C_T.
 \label{eq:logbound}
\end{align}
The minimum in this formula is needed only on admitted pairs.
\end{lemma}
\begin{proof}
For $a,b,c,d\geq0$ and $\delta>0$, put
$L_\delta=\log(a+\delta)-\log(b+\delta)$. Then
\begin{equation}\label{eq:logproduct}
 L_\delta(ac-bd)\geq\min(c,d)L_\delta(a-b)-|a-b||c-d|.
\end{equation}
Indeed, assume $a\geq b$ by interchanging endpoints. If $c\geq d$, write
$ac-bd=d(a-b)+a(c-d)$. If $c<d$, write
$ac-bd=c(a-b)-b(d-c)$ and use
$b\log((a+\delta)/(b+\delta))\leq a-b$.
When $c=a^s,d=b^s$, the negative term in \eqref{eq:logproduct} can be omitted.
Also, Cauchy--Schwarz on the integral of $(t+\delta)^{-1/2}$ gives
\begin{equation}\label{eq:sqrtlog}
 (a-b)L_\delta\geq4|\sqrt{a+\delta}-\sqrt{b+\delta}|^2.
\end{equation}
For any $M>0$ and $0<\delta\leq M$,
\begin{equation}\label{eq:lowhigh}
 |a-b|^2\leq8M|\sqrt{a+\delta}-\sqrt{b+\delta}|^2
                     +M^{-2(s-1)}|a^s-b^s|^2.
\end{equation}
If $\max(a,b)\leq M$, use the square-root factorization. Otherwise, with
$a=\max(a,b)>M$, one has
$(a^s-b^s)/(a-b)\geq a^{s-1}\geq M^{s-1}$.

Let $A=\max_i\sum_{j\ne i}a_{ij}$ and $a_* =\min_i a_{i0}>0$.
For vectors $X,Y$, write $\Delta$ for the difference of their endpoint
values. Summing \eqref{eq:logproduct}--\eqref{eq:sqrtlog}, using symmetry,
and applying Young's inequality to the off-diagonal errors yields
\begin{align*}
 \sum_i\Delta\log(X_i+\delta)\Delta P_i
 \geq{}&4a_*\sum_i|\Delta\sqrt{X_i+\delta}|^2\\
 &+4\sum_{i,j}a_{ij}\min(X_j^s,Y_j^s)|\Delta\sqrt{X_i+\delta}|^2
       -\frac A2\sum_i\bigl(|\Delta X_i|^2+|\Delta X_i^s|^2\bigr).
\end{align*}
For $A>0$ fix $M\leq\min(1,a_*/(2A))$ once and for all. Inserting
\eqref{eq:lowhigh} costs at most $4AM\leq2a_*$ of the first coefficient.
Consequently, for $0<\delta\leq M$,
\begin{align}
 \sum_i\Delta\log(X_i+\delta)\Delta P_i
 \geq{}&2a_*\sum_i|\Delta\sqrt{X_i+\delta}|^2\notag\\
 &+4\sum_{i,j}a_{ij}\min(X_j^s,Y_j^s)|\Delta\sqrt{X_i+\delta}|^2
        -C\sum_i|\Delta X_i^s|^2,\label{eq:logcoercive}
\end{align}
where $C=(A/2)(1+M^{-2(s-1)})$. If $A=0$, take $M=1,C=0$.

Test the fixed-$\eps$ equations with $\log(u_i^\eps+\delta)$, the derivative
of $B_\delta(t)=(t+\delta)\log(t+\delta)-(t+\delta)+1\geq0$.
The chain rule follows as in Lemma~\ref{lem:power}, using a $C^1$ extension
to $t>-\delta/2$ at fixed $\delta$. Uniformly for $\delta\leq1$,
\begin{equation}\label{eq:logreaction}
 \sum_iR_i(u)\log(u_i+\delta)\leq C(1+\sum_i u_i^s).
\end{equation}
For $u_i+\delta\geq1$ discard competition and use
$u_i\log(u_i+1)\leq C_s(1+u_i^s)$. In the other case discard growth and
use $u_i|\log(u_i+\delta)|\leq1/e$, leaving a bound by $C\sum_j u_j$.
The initial $B_\delta$ integral is uniformly bounded. Integrate the entropy
identity, whose diffusion factor is $1/2$, and apply
\eqref{eq:logcoercive}, \eqref{eq:logreaction} and \eqref{eq:powerbound}.
The negative power dissipation has already been bounded. All remaining
dissipation terms are nonnegative. Fatou's lemma as $\delta\downarrow0$
at fixed $\eps$ proves \eqref{eq:logbound}. No positive density floor or
test with an unregularized logarithm at zero is used.
\end{proof}

\section{Compactness and the local equation}\label{sec:limit}

\subsection{Pressure differences and time derivatives}

\begin{lemma}\label{lem:pressure}
For $r=4s/(3s+1)$ and $r'=4s/(s-1)$,
\begin{equation}\label{eq:qpressure}
 \norm{\qe P_i(u^\eps)}_{L^r(\mu)}\leq C_T,\qquad
 \norm{\partial_tu_i^\eps}_{L^r(0,T;(W^{1,r'}(\Om))')}\leq C_T.
\end{equation}
\end{lemma}
\begin{proof}
For $a,b,c,d\geq0$,
\begin{equation}\label{eq:productdiff}
 |ac-bd|\leq\min(c,d)|a-b|+\max(a,b)|c-d|.
\end{equation}
For the product $u_i^\eps(u_j^\eps)^s$, take $a,b$ to be its $i$-species
endpoint values and $c,d$ the $s$th powers of its $j$-species values.
The first term divided by $\eps$ factors, on admitted pairs, as
\[
 \bigl[\sqrt{\min(c,d)}(\sqrt a+\sqrt b)\bigr]
       \bigl[\sqrt{\min(c,d)}\,|\qe\sqrt{u_i^\eps}|\bigr].
\]
For $a_{ij}>0$, the second bracket is bounded in $L^2(\mu)$ by
\eqref{eq:logbound}. The first is bounded in $L^{4s/(s+1)}(\mu)$:
its two factors are in $L^4$ and $L^{4s}$ by \eqref{eq:L2s}.
Endpoint evaluations preserve these norms, by translation of $x$ and the
finite mass of $J$. H\"older's inequality gives their product in $L^r$,
since $1/r=1/2+(s+1)/(4s)$.
The second term of \eqref{eq:productdiff}, divided by $\eps$, is bounded
in $L^{2s/(s+1)}$ by the $L^{2s}$ norm of the endpoint maximum and the
$L^2$ norm of $\qe(u_j^\eps)^s$. The linear pressure term is bounded
in $L^{4s/(2s+1)}$ by factoring its difference through the square root.
Both exponents exceed $r$, and $\mu$ is finite. Terms with $a_{ij}=0$
require no weighted estimate. This proves the first bound.

For $\psi\in W^{1,r'}(\Om)$ the extension theorem on a smooth domain
\cite[Theorem 9.7]{Brezis} and integration along segments give
\begin{equation}\label{eq:testbound}
 \norm{\qe\psi}_{L^{r'}(\Om\times J(z)\dd z)}
                  \leq C\norm{\psi}_{W^{1,r'}(\Om)}.
\end{equation}
Indeed, for an extension $E\psi$ the quotient equals the integral of
$z\cdot\nabla E\psi(x+\theta\eps z)$ over $0\leq\theta\leq1$ on
admitted pairs. Translation and bounded kernel support give the estimate.
Regional symmetry gives
\begin{equation}\label{eq:bilinear}
 \langle\partial_tu_i^\eps,\psi\rangle
 =-\frac12\int_{\Om\times\R^d}J(z)\qe P_i(u^\eps)\qe\psi\dd x\dd z
                       +\int_\Om R_i(u^\eps)\psi.
\end{equation}
Use \eqref{eq:testbound}, the first bound of \eqref{eq:qpressure}, and
the reaction bound in $L^s(Q_T)$. Since $r<s$ and $r'>s'$, the reaction
functional is bounded in the required dual space and time exponent.
\end{proof}

\subsection{Nonlinear space-time compactness}

Fix the negative space $Y=(H_0^m(\Om))'$ with
$m=\lceil d/2\rceil+3$. Sobolev embedding gives
$H_0^m(\Om)\hookrightarrow W^{1,r'}(\Om)\cap L^\infty(\Om)$
\cite[Corollaries 9.13 and 9.15]{Brezis}. Thus $L^1(\Om)$ embeds
continuously and injectively in $Y$; injectivity follows by testing against
$C_c^\infty(\Om)$. Lemma~\ref{lem:pressure} bounds the time derivative in
$L^1(0,T;Y)$.

\begin{lemma}\label{lem:compact}
For each $T<\infty$, the family along any vanishing sequence of $\eps$
is relatively compact in $L^1(Q_T)^n$.
\end{lemma}
\begin{proof}
Consider one species. Choose $C_0$ bounding its spatial $L^s$ norm at all
times by \eqref{eq:powerbound}. If $\eps_k\to0$, $a_k\geq0$,
$\norm{a_k}_s\leq C_0$, and $D_{\eps_k}(a_k^s)\leq M$, then
Lemma~\ref{lem:Ponce} and the bounded mean imply compactness of $a_k^s$
in $L^2$. The roots are compact in $L^{2s}$, since
$|a-b|^s\leq|a^s-b^s|$ for nonnegative $a,b$. In particular they are
compact in $L^1$.

It follows that, for every $\eta>0$, there are $\eps_\eta>0$ and
$C_\eta<\infty$ such that, for $\eps<\eps_\eta$ and
$\norm{a}_s,\norm{b}_s\leq C_0$,
\begin{equation}\label{eq:Ehrling}
 \norm{a-b}_1\leq\eta\bigl[1+\De(a^s)+\De(b^s)\bigr]
                                      +C_\eta\norm{a-b}_Y.
\end{equation}
Here $a,b\geq0$, and infinite energies make the assertion vacuous.
Otherwise, failure would supply $\eps_k<1/k$ violating the estimate with
$C_\eta=k$. The left side has a fixed upper bound, so both energies are
bounded and $\norm{a_k-b_k}_Y\to0$. By the compactness just proved,
subsequences converge in $L^1$. Injectivity of $L^1\hookrightarrow Y$
makes their limits equal. This contradicts the strict lower bound $\eta$
for their $L^1$ distance.

Apply \eqref{eq:Ehrling} to $u_i^\eps(t+h),u_i^\eps(t)$ and integrate.
The integrated $Y$ time shift is at most
$h\norm{\partial_tu_i^\eps}_{L^1(0,T;Y)}$, so
\begin{equation}\label{eq:timeshift}
 \limsup_{\eps\to0}\int_0^{T-h}
       \norm{u_i^\eps(t+h)-u_i^\eps(t)}_1\dd t
           \leq\eta(T+2C_T)+C_\eta C_T h.
\end{equation}
First send $h\to0$, then $\eta\to0$.

For spatial translations, extend each density by zero outside $\Om$.
On slices with $\De((u_i^\eps)^s)\leq M$, the spatial compactness above
gives an asymptotically uniform $L^1$ translation modulus. To see this,
a failure would give $\eps_k\to0$, shifts $h_k\to0$ and such slices
converging in $L^1$ but with a nonvanishing translation defect. This
contradicts translation continuity of the fixed zero-extended limit.
The complementary times have measure at most $C_T/M$ by
\eqref{eq:powerbound}. Their translation contribution is at most
$2\sup_t\norm{u_i^\eps(t)}_1 C_T/M$. Send the spatial shift to zero
on the good slices, then $M\to\infty$.

After zero extension in time, the endpoint strips contribute at most $C_T h$.
All extensions have common bounded support in $\R^{d+1}$. For any fixed
vanishing sequence, finitely many early members have the ordinary translation
continuity of a single $L^1$ function. Combining them with the eventual
estimates and the triangle inequality gives uniform space-time translation
continuity. The Fr\'echet--Kolmogorov theorem
\cite[Theorem 4.26 and Lemma 4.3]{Brezis} proves the assertion.
\end{proof}

After extraction, $u_i^\eps\to u_i$ strongly in $L^1(Q_T)$ and almost
everywhere. Fatou's lemma and \eqref{eq:L2s}, followed by interpolation,
give \eqref{eq:conv}. For $p=s+1<2s$ the elementary estimate
\[
 |a^s-b^s|\leq s|a-b|(a+b)^{s-1}\quad(a,b\geq0)
\]
gives strong convergence of $(u_j^\eps)^s$ in $L^{p/s}$.
Splitting the product difference and using H\"older with $p$ and $p/s$
then gives convergence of $u_i^\eps(u_j^\eps)^s$ in $L^1$.
Strong $L^2$ convergence similarly passes all quadratic reaction products.
This proves \eqref{eq:productconv}, without a weak-product argument.
Nonnegativity and the $L^\infty_tL^s_x$ and $L^{2s}$ bounds pass to the limit.

\subsection{The pressure gradient, boundary condition and initial trace}

\begin{proof}[Proof of Theorem~\ref{thm:main}]
We first work on a fixed finite horizon. Recall the regional quotient and its
fixed measure:
\[
 \qe w=\mathbf1_\Om(x+\eps z)
       \frac{w(t,x+\eps z)-w(t,x)}{\eps},\qquad
 \dd\mu=\dd t\dd x\,J(z)\dd z.
\]
By \eqref{eq:qpressure} and reflexivity, a further subsequence satisfies
$\qe P_i(u^\eps)\rightharpoonup G_i$ in $L^r(\mu)$, where $r>1$.
Take $\eta\in C_c^\infty((0,T)\times\Om)$ and a bounded measurable
function $\zeta(z)$ on $\supp J$. Translating the first term of the
quotient, for small enough $\eps$ relative to the distance of $\supp\eta$
from the boundary, gives in the limit
\begin{equation}\label{eq:G}
 \int G_i\eta\zeta\dd\mu
       =-\int P_i(u)\,z\cdot\nabla_x\eta\,\zeta\dd\mu.
\end{equation}
Here \eqref{eq:productconv} pairs with the uniformly convergent smooth test
difference quotient. We have not applied a difference operator to a bare
$L^1$ limit.
Define $V_{ik}(t,x)=\frac12\int J(z)z_kG_i(t,x,z)\dd z$.
H\"older's inequality and compact kernel support give $V_{ik}\in L^r(Q_T)$.
Use $\zeta(z)=z_k$ in \eqref{eq:G} and \eqref{eq:moment}; then
$V_{ik}=\partial_kP_i(u)$ in distributions. The pressure itself belongs
to $L^r$ by its $L^{2s/(s+1)}$ bound. Hence
$P_i(u)\in L^r(0,T;W^{1,r}(\Om))$.
Returning to all product tests in \eqref{eq:G}, distributional separation
and Fubini give
\begin{equation}\label{eq:qid}
 G_i(t,x,z)=z\cdot\nabla P_i(u)(t,x)\quad\mu\text{-almost everywhere}.
\end{equation}
The identical argument at exponent two, using the power and square-root
difference bounds, proves $u_i^s,\sqrt{u_i}\in L^2(0,T;H^1(\Om))$.
Their $L^2$ norms follow from \eqref{eq:L2s} and the mass bound, respectively.

Now let $\phi\in C^1([0,T]\times\overline\Om)$, $\phi(T)=0$.
Extend the spatial test by the bounded $W^{1,\infty}$ extension theorem.
Its Lipschitz representative bounds $|\qe\phi|$ by $C|z|$. At every
interior point $x$, for small enough $\eps$, the quotient uses only the
original $C^1$ function. The regional indicator tends to one almost
everywhere. Dominated convergence therefore gives
\[
 \qe\phi\longrightarrow z\cdot\nabla\phi
             \quad\hbox{strongly in }L^{r'}(\mu).
\]
Integrate \eqref{eq:bilinear} in time and use
\eqref{eq:productconv} and \eqref{eq:qid}. The diffusion pairing tends to
\[
 \frac12\int J(z)(z\cdot\nabla P_i(u))(z\cdot\nabla\phi)\dd z
                         =\nabla P_i(u)\cdot\nabla\phi.
\]
The initial term remains exactly $\int_\Om u_i^0\phi(0)$.
This proves \eqref{eq:weak} for arbitrary boundary values of $\phi$.
Its absence of a boundary contribution is precisely the weak homogeneous
Neumann pressure condition.

For completeness we identify the time representative, including at zero.
Put $V=W^{1,r'}(\Om)$ and define
\[
 F_i(t)(\psi)=-\int_\Om\nabla P_i(u)\cdot\nabla\psi
                         +\int_\Om R_i(u)\psi.
\]
Its $V'$ norm is bounded by an $L^r(0,T)$ function. Smooth functions up to
the boundary are dense in $V$: extend to $\R^d$, mollify, use translation
continuity for the function and its weak derivatives, then restrict to $\Om$.
This uses the extension theorem and translation lemma already cited, and
does not assert density of interior compact tests in the full boundary space.
Thus \eqref{eq:weak} shows that $u_i$ agrees almost everywhere with
\[
 U_i(t)=u_i^0+\int_0^tF_i(\tau)\dd\tau\quad\hbox{in }V'.
\]
This representative is continuous in $V'$. At an exceptional time,
approach by ordinary times, use the uniform $L^s$ bound and reflexivity to
obtain a weak $L^s$ limit. Continuity in $V'$ identifies it uniquely with
$U_i(t)$, so the bound holds at every time. Pairings with arbitrary
$L^{s'}$ tests are uniform limits of continuous pairings with smooth tests,
by their density in $L^{s'}$ and the uniform $L^s$ bound.
Therefore $u_i\in C_w([0,T];L^s)$, $u_i(0)=u_i^0$ weakly in $L^s$,
and its time derivative has the asserted $L^r(0,T;V')$ bound.

Finally, start with any vanishing sequence and perform nested extractions
on $T=1,2,3,\ldots$. The diagonal subsequence converges on every finite
horizon. Its strong $L^1$ limits agree on overlaps, defining one global
$u$. All estimates and weak identities restrict to any finite $T$.
This proves the theorem.
\end{proof}

\begin{thebibliography}{9}
\bibitem{AIM} AIM open-problem catalogue, Problem 469,
\emph{The local limit of superlinear multispecies nonlocal diffusion}, revision
\texttt{a153f469ee0ea64e876b2ad97c79d252c19a3e22}.
\href{https://github.com/MColbrook/AIM/blob/a153f469ee0ea64e876b2ad97c79d252c19a3e22/problems/469-superlinear-skt-nonlocal-to-local-limit.md}{Original problem statement}.
\bibitem{Brezis} H. Brezis, \emph{Functional Analysis, Sobolev Spaces and
Partial Differential Equations}, Springer, 2011.
\href{https://doi.org/10.1007/978-0-387-70914-7}{doi:10.1007/978-0-387-70914-7}.
\bibitem{CDJ} X. Chen, E. S. Daus and A. J\"ungel,
Global existence analysis of cross-diffusion population systems for multiple
species, \emph{Arch. Ration. Mech. Anal.} \textbf{227} (2018), 715--747.
\href{https://doi.org/10.1007/s00205-017-1172-6}{doi:10.1007/s00205-017-1172-6}.
\bibitem{GV} G. Galiano and J. Velasco, Convergence of solutions of a
rescaled evolution nonlocal cross-diffusion problem to its local diffusion
counterpart, \emph{Rev. R. Acad. Cienc. Exactas F\'is. Nat. Ser. A Mat.
RACSAM} \textbf{116} (2022), article 93.
\href{https://doi.org/10.1007/s13398-022-01231-7}{doi:10.1007/s13398-022-01231-7}.
\bibitem{HJP} P. Hirvonen, A. J\"ungel and A. Pollino,
Superlinear nonlocal diffusion systems for multispecies populations:
well-posedness and discrete chain rules,
\href{https://arxiv.org/abs/2609.09921v1}{arXiv:2609.09921v1}, 2026.
\bibitem{Ponce} A. C. Ponce, An estimate in the spirit of Poincar\'e's
inequality, \emph{J. Eur. Math. Soc.} \textbf{6} (2004), 1--15.
\href{https://doi.org/10.4171/JEMS/1}{doi:10.4171/JEMS/1}.
\end{thebibliography}

\section*{AI-use disclosure}
The mathematical argument and manuscript were developed with assistance from
OpenAI Codex, including AI-assisted proof checks and manuscript reviews.
These checks do not constitute independent human review or formal verification.
No numerical computation is a premise of the result.
\end{document}
